\section{Cahier des charges}
\redacteur{\acompleter[nom]}

\consigne{Présentez le MVP tel que validé en semaine 1, puis ce qui a réellement été livré. Tout écart doit être expliqué : un écart justifié n'est pas pénalisé, un écart caché l'est (critère P1 de la grille).}

\subsection{Fonctionnalités du MVP}

\begin{table}[H]\centering
\caption{Fonctionnalités prévues et état à la livraison}\label{tab:mvp}
\begin{tabularx}{\linewidth}{lXll}
\toprule
\textbf{N°} & \textbf{Fonctionnalité} & \textbf{Prévue en S1} & \textbf{État final} \\
\midrule
F1 & \acompleter & oui & livrée et testée \\
F2 & \acompleter & oui & \acompleter \\
F3 & \acompleter & oui & \acompleter \\
\bottomrule
\end{tabularx}
\end{table}

\subsection{Exigence liée à la taille de l'équipe}
\consigne{Trinôme : composant supplémentaire (évaluation automatisée, monitoring ou second agent). Équipe de 4 : deux agents ou deux sous-systèmes intégrés. Binôme : écrire « sans objet ».}
\acompleter

\subsection{Hors périmètre et écarts avec la semaine 1}
\acompleter
